% TÀI LIỆU / BÀI TẬP CẤP 2: tiêu đề, phần, dạng, đề mục loại bài.
% Không có danh mục nên đặt tên tài liệu bằng \dattentailieu (dùng \loai, \chuong thì phải có danh mục).
\documentclass[12pt]{article}
\usepackage{toan_thcs}
\begin{document}
\mon{TOÁN 9}
\dattentailieu{Chuyên đề Căn bậc hai}
\tieude{\tentailieu}{Bài 1. Căn bậc hai số học}

\phan{A. Lý thuyết}
\begin{kienthuc}
Căn bậc hai số học của số $a\ge0$ là số $x\ge0$ sao cho $x^2=a$.
\end{kienthuc}

\phan{B. Phân loại và phương pháp giải toán}
\dang{Tính giá trị căn thức}
\muctl
\begin{bt}[TH]
Tính $\sqrt{49}+\sqrt{16}$.
\loigiai{$7+4=11$.}
\end{bt}
\muctn
\begin{tn}[NB]
Giá trị của $\sqrt{25}$ là
pa{$4$.}{dung{$5$.}}{$-5$.}{$25$.}
\end{tn}
\mucds
\begin{ds}
Cho $a=\sqrt{36}$.
\ydung{$a=6$.}
\ysai{$a=-6$.}
\ydung{$a^2=36$.}
\ysai{$a=18$.}
\end{ds}
\muctln
\begin{tln}
Tính $\sqrt{121}$.
\traloi{11}
\end{tln}
\bangdapan
\end{document}
